\section{Actualización del sistema de benchmarks: de los benchmarks web a OS-World}

\subsection{El valor y los límites de los benchmarks existentes}
El expositor repasó benchmarks como Mind2Web, WebArena y VisualWebArena. Estos trabajos impulsaron la investigación en Web Agent, pero su cobertura de entornos, tipos de tarea y escalabilidad siguen sin bastar para representar escenarios reales de ``computer use'', en particular las tareas entre aplicaciones y de objetivo abierto.

\begin{knowledgebox}{Limitaciones típicas}
\begin{itemize}
    \item La distribución de tareas se inclina hacia la navegación web, con una cobertura estrecha;
    \item Muchas tareas siguen rutas implícitamente plantilladas, fáciles de sobreajustar con una política;
    \item El soporte para aplicaciones de escritorio, sistemas de archivos y flujos entre programas es insuficiente;
    \item Es difícil sostener un ciclo integrado de ``desarrollo, prueba y evaluación''.
\end{itemize}
\end{knowledgebox}

\subsection{Los objetivos de diseño de OS-World}
OS-World se posiciona como un entorno informático real más general: dentro de una máquina virtual ofrece un escritorio interactivo y un ecosistema de aplicaciones, y permite que los investigadores definan tareas, ejecuten agentes, recolecten trayectorias y corran scripts de evaluación, dejando además espacio para ampliarlo.

\begin{importantbox}{Tres propiedades clave de OS-World}
\begin{enumerate}
    \item \textbf{Realismo}: se acerca al comportamiento real de un sistema operativo y sus aplicaciones;
    \item \textbf{Escalabilidad}: permite agregar de manera sostenida nuevas tareas y combinaciones de aplicaciones;
    \item \textbf{Evaluabilidad}: cada tarea puede vincularse a una lógica de verificación programada.
\end{enumerate}
\end{importantbox}

\begin{figure}[H]
\centering
\includegraphics[width=0.88\textwidth]{frame-02-osworld.jpg}
\caption{Video aproximadamente en 00:10:32: se propone OS-World como un entorno real y escalable para el desarrollo y la evaluación de agentes multimodales}
\end{figure}

\subsection{Perspectiva comparativa: por qué OS-World se acerca más a los problemas de producción}
\begin{table}[H]
\centering
\begin{tabular}{p{3.2cm}p{3.7cm}p{3.8cm}}
\toprule
\textbf{Dimensión} & \textbf{Benchmarks web tradicionales} & \textbf{OS-World} \\
\midrule
Cobertura de entornos & Principalmente el dominio web & Web + aplicaciones de escritorio + interacción con el sistema \\
Forma de las tareas & Se inclina hacia navegación y formularios & Flujos entre aplicaciones, de varios pasos y de objetivo abierto \\
Método de evaluación & Orientado al resultado de la página & Verificación por script + recompensa por cumplimiento parcial \\
Capacidad de ampliación & Costo alto para construir tareas nuevas & Soporta agregar tareas y configuraciones de manera sistemática \\
Adaptación de ingeniería & Se inclina hacia la evaluación fuera de línea & Desarrollo, entrenamiento y evaluación integrados \\
\bottomrule
\end{tabular}
\caption{OS-World no busca sustituir todos los benchmarks, sino llenar el vacío central de las ``tareas informáticas reales''}
\end{table}

\begin{warningbox}{Riesgo de traslado de la evaluación}
Si un equipo solo optimiza sobre un único benchmark, el modelo tiende a desarrollar una ``especulación de política'' frente al script de evaluación. El curso recomienda hacer validación cruzada en múltiples entornos y familias de tareas para evitar el sobreajuste al benchmark.
\end{warningbox}

\subsection{Resumen de la sección}
El significado de OS-World no es simplemente sumar un benchmark más, sino desplazar la investigación de agentes de la ``competencia de bancos de preguntas'' hacia la ``ingeniería de entornos''. Esto ofrece una base experimental unificada para el diseño de datos y modelos posterior.
